% =========================================================================
% Apendice B - Construcao das variaveis e codigos
% =========================================================================

\section{Construção das variáveis e códigos}
\label{app:variaveis}

Conforme discutido na Seção~\ref{sec:dados_metodo}, algumas variáveis
educacionais da PNADC mudaram de nome e de codificação ao longo do
período 2012--2024. A Tabela~\ref{tab:harmonizacao} na seção principal
lista os nomes das variáveis nas duas versões. A
Tabela~\ref{tab:codigos_v3003} a seguir detalha os mapeamentos de
valores entre as duas codificações para as duas variáveis mais
críticas, \texttt{V3003} pré-2016 e \texttt{V3003A} a partir de 2016.

\begin{table}[htbp]
\centering
\caption{Mapeamento de códigos: \texttt{V3003} (2012--2015) $\to$ \texttt{V3003A} (2016+).}
\label{tab:codigos_v3003}
\small
\begin{tabular}{cclc}
\toprule
\texttt{V3003} & \texttt{V3003A} & Curso/etapa frequentado & Macroetapa (esta nota) \\
\midrule
1 & 1 & Creche                                          & ---             \\
2 & 2 & Pré-escola                                      & ---             \\
3 & 4 & Regular Fundamental                             & EF iniciais / finais \\
4 & 5 & EJA Fundamental                                 & EJA EF          \\
5 & 6 & Regular Médio                                   & EM              \\
6 & 7 & EJA Médio                                       & EJA EM          \\
7 & 8 & Educação Profissional (médio)                   & EM (técnico)    \\
8 & 10 & Superior - graduação                           & ---             \\
9 & 11 & Especialização (nível superior)                & ---             \\
10 & 12 & Mestrado                                      & ---             \\
11 & 13 & Doutorado                                     & ---             \\
--- & 3  & Alfabetização de jovens e adultos            & EJA EF          \\
--- & 9  & Educação Profissional (qualif.\ profissional) & ---             \\
\bottomrule
\end{tabular}
\end{table}

A discrepância crítica é o deslocamento de uma unidade para os códigos
maiores ou iguais a três a partir de 2016, decorrente da inserção de
novas categorias, especificamente da ``Alfabetização de jovens e
adultos'', código 3 na nova versão. Sem a harmonização, alunos
classificados como ``Regular Fundamental'' em 2012--2015 (código 3
em \texttt{V3003}) seriam erroneamente classificados como
``Alfabetização'' (código 3 em \texttt{V3003A}) caso as variáveis
fossem empilhadas sem ajuste, erro que afeta cerca de noventa e cinco
por cento das observações pré-2016 em idade escolar EF. A variável de
série, \texttt{V3006}, é estável em todo o período 2012--2024.
Codifica o ano que o indivíduo cursa, no formato $XYY$, em que $X$ é
o nível (1 para Fundamental, 2 para Médio, 3 para Superior e 4 para
outros) e $YY$ é o ano (01 a 09 para EF e 01 a 04 para EM). Por
exemplo, \texttt{V3006=108} identifica o oitavo ano do EF. A
variável \texttt{V3006A}, introduzida em 2016, acrescenta a
desagregação ``EF anos iniciais'' (1 a 5) e ``EF anos finais'' (6 a
9), redundante com \texttt{V3006} mas conveniente para filtros.

O identificador domiciliar estável é construído como a concatenação
$\texttt{hh\_id} = \texttt{UF} \,\|\, \texttt{UPA} \,\|\, \texttt{V1008} \,\|\, \texttt{V1014}$,
e o identificador individual pré-validação é
$\texttt{pid} = \texttt{hh\_id} \,\|\, \texttt{V2003}$. Após a
validação de consistência por sexo e idade, elos com inconsistência
em vinte por cento ou mais das visitas são marcados como elo
quebrado e descartados do painel. Para corrigir possível viés de
atrito, estimamos a propensão de retenção como
\[
\Pr(R_i = 1 \mid X_i) = \Phi(X_i' \beta),
\]
em que $R_i \in \{0,1\}$ indica retenção do indivíduo $i$ no painel
entre as duas visitas usadas no indicador, e $X_i$ inclui idade,
idade ao quadrado, sexo, raça, renda per capita, logaritmo da renda
per capita, rede da escola, UF (efeito fixo) e trimestre da primeira
observação (efeito fixo). O peso corrigido por \emph{inverse
propensity}, denominado P3, é
\[
w_i^{P3} = w_i^{P1} \cdot \frac{1}{\hat\Pr(R_i=1 \mid X_i)}.
\]

Toda a análise é reprodutível usando o código disponível em
\href{https://github.com/vitpereira/Nota_PNAD}{github.com/vitpereira/Nota\_PNAD}
(a ser publicado na submissão). O \emph{stack} computacional é
composto por Python para o \emph{parse} dos microdados, Stata 19
para a construção do painel e o cálculo dos indicadores e R com
\emph{ggplot2} para as figuras. Sementes aleatórias estão fixadas
para garantir a reprodutibilidade dos testes Monte Carlo. Para
facilitar a replicação e a aplicação desta metodologia em outros
contextos, disponibilizamos um pacote R público em
\href{https://github.com/vitpereira/fluxoescolar}{github.com/vitpereira/fluxoescolar}.
O pacote oferece oito funções exportadas, a saber,
\texttt{baixar\_pnadc()} para download direto do FTP do IBGE,
\texttt{parsear\_pnadc()} para \emph{parse} dos microdados
fixed-width, \texttt{harmonizar\_pnadc()} para consolidação dos
códigos pré e pós 2016, \texttt{construir\_painel()} para linkagem
individual via Ribas-Soares, \texttt{calcular\_indicadores()} para
cálculo dos cinco indicadores, \texttt{decompor\_inep()} para a
decomposição $R+U+S+C+M$, e duas funções de visualização. O pacote
inclui testes unitários, \emph{vignette} introdutória e fluxo de
CI via GitHub Actions, e pode ser instalado via\\
\texttt{remotes::install\_github("vitpereira/fluxoescolar")}.
